\section{Results}

\subsection{Research question 1 (RQ1): Is Tissue-specific Preference Learnable?}

We first ask whether peptide sequence contains any useful tissue-associated signal after the \plainMHC{} restriction and source protein are held constant. All current human results use the same eligibility rule as the mouse benchmark: each retained tissue--antigen-presenting-molecule combination (task) contains strictly more than 200 eligible peptide pairs sharing restriction and parent protein (matched pairs). On the resulting ordinary human benchmark with 157 biological combinations (157-task benchmark), even the model with one peptide representation shared by all combinations (\plainSharedHeads{}) reaches an AUROC averaged equally across combinations (mean task AUROC) of 0.7972 and AUPRC of 0.7829 (Table~\ref{tab:human-standard}), well above random ordering. Models with more structured information sharing increase AUROC to 0.8239--0.8448. This shows that peptides matched by presenting restriction and parent protein are not exchangeable at random: their sequences contain patterns associated with the recorded tissue context. It does not show whether those patterns arise from antigen processing, protein abundance, cell composition, or study structure.

\plainPeptideOOF{} performance is reported once, together with its \plainMatchedOOF{} reference and uncertainty, under research question 4 (RQ4; Tables~\ref{tab:human-standard-strict-summary} and~\ref{tab:mouse-standard-strict-summary}); this avoids repeating the same result across research questions 1--3 (RQ1--RQ3).

\subsection{Research question 2 (RQ2): Main Human Benchmark Performance}

On the ordinary human benchmark with 157 tissue--\plainHLA{} combinations (157-task benchmark), TissuePMHC reaches AUROC/AUPRC values averaged equally across combinations (mean task AUROC/AUPRC) of 0.8448/0.8348. In practical terms, peptides reported in the target tissue usually receive higher scores than comparison peptides sharing restriction and parent protein (matched comparisons), although performance varies markedly across tissue--\plainHLA{} combinations. Accuracy is 0.7772, MCC is 0.5545, and the mean AUROC of the ten most difficult combinations (tasks) falls to 0.6834. Table~\ref{tab:human-standard} compares models on the same set of tissue--\plainHLA{} combinations (task inventory).

\begin{table}[tbp]
\centering
\caption{Human \plainFixedTest{} benchmark across 157 tissue--\plainHLA{} tasks. Rows labelled as a mean are arithmetic summaries of three separately trained model results (member-level metric tables); rows labelled as an average across models (ensemble) are calculated after averaging row-level predictions. Only the two averaged-model rows use directly comparable aggregation.}
\label{tab:human-standard}
\footnotesize
\setlength{\tabcolsep}{2pt}
\begin{tabularx}{\textwidth}{>{\raggedright\arraybackslash}Xrrrrrr}
\toprule
Model & \shortstack{Mean\\AUROC} & \shortstack{Median\\AUROC} & AUPRC & Acc. & MCC & Worst-10 \\ 
\midrule
\plainSharedHeads{} (mean across three models) & 0.7972 & 0.8068 & 0.7829 & 0.7277 & 0.4598 & 0.6299 \\ 
\plainAuxDual{} (mean across three models) & 0.8239 & 0.8351 & 0.8117 & 0.7552 & 0.5132 & 0.6501 \\ 
Two-view \plainMLP{} (averaged model; ensemble) & 0.8345 & 0.8441 & 0.8233 & 0.7655 & 0.5313 & 0.6666 \\ 
TissuePMHC (averaged model; ensemble) & \textbf{0.8448} & \textbf{0.8585} & \textbf{0.8348} & \textbf{0.7772} & \textbf{0.5545} & \textbf{0.6834} \\ 
\bottomrule
\end{tabularx}
\end{table}

\begin{figure}[tbp]
\centering
\begin{tikzpicture}
\begin{axis}[
    width=0.97\linewidth,
    height=6.4cm,
    boxplot/draw direction=y,
    ymin=0.58, ymax=1.0,
    xmin=0.4, xmax=4.6,
    xtick={1,2,3,4},
    xticklabels={Shared representation,Training guidance,Simple two-view model,TissuePMHC},
    x tick label style={rotate=18,anchor=east,font=\small},
    ylabel={AUROC for each tissue--HLA combination (task-level AUROC)},
    grid=major, grid style={gray!20}
]
\addplot+[boxplot prepared={draw position=1,lower whisker=0.6066,lower quartile=0.7305,median=0.8009,upper quartile=0.8681,upper whisker=0.9722},fill=blue!12,draw=blue!65!black] coordinates {};
\addplot+[boxplot prepared={draw position=2,lower whisker=0.6314,lower quartile=0.7438,median=0.8305,upper quartile=0.9033,upper whisker=0.9835},fill=blue!18,draw=blue!65!black] coordinates {};
\addplot+[boxplot prepared={draw position=3,lower whisker=0.6424,lower quartile=0.7582,median=0.8441,upper quartile=0.9107,upper whisker=0.9859},fill=blue!24,draw=blue!65!black] coordinates {};
\addplot+[boxplot prepared={draw position=4,lower whisker=0.6510,lower quartile=0.7632,median=0.8585,upper quartile=0.9287,upper whisker=0.9878},fill=blue!30,draw=blue!65!black] coordinates {};
\addplot+[only marks,mark=*,mark size=2.5pt,red!75!black] coordinates {(1,0.7972) (2,0.8239) (3,0.8345) (4,0.8448)};
\end{axis}
\end{tikzpicture}
\caption{Human ordinary-benchmark distributions of AUROC calculated separately for each tissue--\plainHLA{} combination (task-level AUROC). Boxes show the first quartile, median, and third quartile; whiskers show the observed minimum and maximum; red dots show means across biological combinations (task means). The first two distributions summarize three separately trained model tables, whereas the last two use predictions averaged across models (row-level prediction ensembles), matching Table~\ref{tab:human-standard}.}
\label{fig:human-standard}
\end{figure}

Relative to the directly comparable \plainMLP{} averaged across three \plainSeeds{}, TissuePMHC improves mean AUROC by 0.0103 and the difficult-task summary by 0.0168. The improvement is visible in the mean, median, and lower-performing tasks, rather than being driven only by a few easy tissue--\plainHLA{} combinations. Larger differences from the \plainSharedHeads{} are descriptive because the aggregation procedures are not identical.

We also repeat the comparison using \plainPairOOF{}, which rotates complete held-out pairs through the training pool. This check asks whether the advantage appears beyond one particular \plainFixedTest{}. The closest comparison model uses the same two-pathway design (two-view training design) but a simpler peptide representation, making it a more informative reference than a model that shares everything globally.

\begin{table}[tbp]
\centering
\caption{\plainPairOOF{} comparison on the 157-task human benchmark. Both models use the same folds; this is distinct from the \plainMatchedOOF{} reference used specifically for analysis of peptide separation across all tasks.}
\label{tab:human-matched-oof}
\small
\begin{tabularx}{\textwidth}{>{\raggedright\arraybackslash}Xrrr}
\toprule
Model & AUROC & AUPRC & Worst-10 \\ 
\midrule
\plainAuxDual{} on identical folds & 0.8131 & 0.7980 & 0.6419 \\ 
TissuePMHC & \textbf{0.8279} & \textbf{0.8144} & \textbf{0.6562} \\ 
\bottomrule
\end{tabularx}
\end{table}

The \plainPeptideOOF{} result is not repeated here; research question 4 (RQ4) reports it with the \plainMatchedOOF{} reference, a scatter plot comparing the same biological combinations (task-wise scatter), an uncertainty interval, and a complete four-metric comparison.

\subsection{Research question 3 (RQ3): Component Contributions on the Human Benchmark}

Table~\ref{tab:component-evidence} summarizes three component-level questions on the current 157-task human inventory: whether preserving residue positions while examining several motif lengths improves performance, whether sharing across all tasks and sharing within one \plainHLA{} allele provide complementary rankings, and whether averaging several \plainSeeds{} improves stability. These are completed system comparisons rather than a design in which every component is changed alone while all others remain fixed (full factorial ablation). Replacing the simpler peptide representation with a position-aware encoder that examines several motif lengths (\plainMultiKernel{}) raises mean AUROC from 0.8345 to 0.8448 and the difficult-task summary from 0.6666 to 0.6834.

The pathway trained across all tasks performs better on average than the pathway restricted to one \plainHLA{} allele, but combining their within-task orderings (rank fusion) reaches 0.8225--0.8249 per run, above either pathway alone. The allele-restricted pathway (HLA-specific view) therefore contributes useful peptide ordering for some tasks even though its overall average is lower.

Third, TissuePMHC trained from three separate \plainSeeds{} obtains mean task AUROCs of 0.8345, 0.8314, and 0.8330 (mean 0.8329; sample standard deviation 0.0016), whereas predictions averaged across the three \plainSeeds{} (three-seed ensemble) reach 0.8448. Averaging \plainSeeds{} therefore provides a measurable reduction in random training variation.

\begin{table}[tbp]
\centering
\caption{Evidence about system components from completed runs on the current 157-task human benchmark. This is not one experiment that changes each component alone (factorial ablation), and no row should be interpreted as an isolated component effect under \plainPeptideOOF{}.}
\label{tab:component-evidence}
\scriptsize
\setlength{\tabcolsep}{2pt}
\begin{tabular}{@{}>{\raggedright\arraybackslash}p{2.3cm}>{\raggedright\arraybackslash}p{4.0cm}>{\centering\arraybackslash}p{3.0cm}>{\centering\arraybackslash}p{1.4cm}@{}}
\toprule
Design question & Completed comparison & AUROC & Descriptive $\Delta$ \\ 
\midrule
Position-aware filters of several widths (multi-kernel encoder) & \plainMLP{} averaged over three \plainSeeds{} $\rightarrow$ TissuePMHC (157 tasks) & 0.8345 $\rightarrow$ 0.8448 & +0.0103 \\ 
Complementarity of the two pathways (branch complementarity) & Better single pathway $\rightarrow$ combined within-task orderings (rank fusion; 157 tasks; ranges across \plainSeeds{}) & 0.8072--0.8091 $\rightarrow$ 0.8225--0.8249 & +0.0134--0.0177 \\ 
Averaging random initializations (seed averaging) & Mean of separate \plainSeeds{} $\rightarrow$ prediction average across three \plainSeeds{} (ensemble; 157 tasks) & 0.8329 $\rightarrow$ 0.8448 & +0.0118 \\ 
\bottomrule
\end{tabular}
\end{table}

Taken together, the completed ordinary-benchmark runs provide descriptive evidence for the assembled system design, but they do not constitute a uniform causal decomposition of all components. The controls evaluated on identical \plainPeptideOOF{} folds under research question 7 (RQ7) provide the stronger comparison: additional tissue and allele guidance during training (auxiliary supervision) retains a clear benefit after peptide separation, whereas the \plainPlainDual{} with combined within-task orderings and the \plainMultiKernel{} relative to the strongest \plainAuxDual{} using an \plainMLP{} do not show a stable AUROC advantage.

\subsection{Research question 4 (RQ4): Effect of Evaluating Peptides Absent from All Training Tasks}

\plainMatchedOOF{} is created by repartitioning the complete training pair pool, not by selecting a subset of an existing test set. Within each tissue--antigen-presenting-molecule combination (task), complete pair identifiers are shuffled with the fixed split seed and assigned round-robin to three folds; for each fold, the other two folds form a newly reconstructed fitting set and the assigned fold forms the held-out set. This keeps the two rows of every pair together but places no constraint on peptide identity: a held-out peptide may reappear in the fitting rows of the same task or of another task. \plainPeptideOOF{} starts from the same complete pair pool but assigns entire connected components of pairs linked by shared peptide sequences to folds, so every held-out peptide is absent from every fitting task. Both protocols retain previously represented tissue--antigen-presenting-molecule combinations (tissue--MHC combinations), use the same pooled evaluation rows, and have nearly identical held-out sizes for each biological combination (task-wise held-out sizes). Their difference describes the change in performance under exact-peptide separation within known biological contexts; it does not test an unseen tissue, allele, task, or source protein, and it is not a pure causal estimate because the stricter grouping can also change fold composition.

\subsubsection{Human benchmark}

Removing peptide reuse makes the human biological ranking problem (task) clearly harder (Table~\ref{tab:human-standard-strict-summary}). The ordinary and peptide-separated analyses evaluate the same 73,798 human pairs. Across the 471 task-by-fold cells, 333 have exactly the same held-out pair count; the remaining cells differ by one pair, giving a mean absolute difference of 0.293 pair. In the ordinary folds, 69.19--69.91\% of unique held-out peptides reappear in fitting, compared with 0\% after global peptide separation. Thus, ``matched'' denotes matching the pair pool, model family, seeds, fold count, pooled evaluation rows, and nearly all task-by-fold sizes; it does not mean retaining an original test set or forcing corresponding pairs into the same fold.

\begin{table}[tbp]
\centering
\caption{Human AUROC under \plainMatchedOOF{} and \plainPeptideOOF{}. Difference is peptide-separated minus ordinary same-pool cross-validation; a negative value denotes lower performance after exact peptide separation.}
\label{tab:human-standard-strict-summary}
\small
\begin{tabularx}{\textwidth}{>{\raggedright\arraybackslash}X
                            >{\raggedright\arraybackslash}Xr}
\toprule
\plainMatchedOOF{} & \plainPeptideOOF{} & Difference \\ 
\midrule
0.8280 & 0.7652 & -0.0628 \\ 
\bottomrule
\end{tabularx}
\end{table}

The relationship across human tissue--\plainHLA{} combinations (task-wise relationship) and the equal-performance reference are shown in Figure~\ref{fig:human-standard-strict-scatter}.
\begin{figure}[tbp]
\centering
\begin{tikzpicture}
\begin{groupplot}[
    group style={group size=1 by 1},
    width=0.72\linewidth,
    height=0.52\linewidth,
    xmin=0.5,xmax=1.0,
    ymin=0.5,ymax=1.0,
    xlabel={Same-pool ordinary cross-validation (matched standard OOF) AUROC},
    ylabel={Peptide-separated cross-validation (global peptide-disjoint OOF) AUROC},
    grid=major,
    grid style={gray!20}
]
\nextgroupplot[title={Human tissue--HLA combinations}]
\addplot+[only marks,mark=*,mark size=1.2pt,blue!70] coordinates {(0.96672,0.89632) (0.97035,0.81086) (0.70311,0.67863) (0.80304,0.78235) (0.73140,0.70095) (0.68931,0.62863) (0.77000,0.71465) (0.62590,0.62946) (0.74105,0.70654) (0.75984,0.69934) (0.77139,0.71416) (0.72978,0.72981) (0.75635,0.73227) (0.77172,0.75377) (0.65427,0.63011) (0.71653,0.68030) (0.80139,0.74274) (0.78455,0.76390) (0.77786,0.75187) (0.81349,0.78630) (0.74166,0.76571) (0.82981,0.78587) (0.73312,0.72855) (0.83455,0.83125) (0.85514,0.83040) (0.95176,0.86541) (0.90220,0.70352) (0.92871,0.75592) (0.87940,0.78940) (0.91259,0.87720) (0.88752,0.84003) (0.70758,0.65148) (0.74956,0.74037) (0.77205,0.76026) (0.74550,0.72139) (0.91297,0.87429) (0.91673,0.84768) (0.86843,0.80669) (0.91822,0.86259) (0.92345,0.75413) (0.87903,0.78704) (0.92050,0.80466) (0.86798,0.78964) (0.94395,0.72269) (0.91009,0.75161) (0.85547,0.74550) (0.82809,0.64234) (0.96895,0.89676) (0.91098,0.71989) (0.87057,0.80826) (0.93357,0.81315) (0.94807,0.77722) (0.82663,0.72979) (0.91800,0.83660) (0.95452,0.87204) (0.84433,0.80015) (0.92370,0.83370) (0.96999,0.80102) (0.78527,0.66459) (0.87280,0.84343) (0.88053,0.80683) (0.66490,0.65294) (0.77697,0.66023) (0.86138,0.78989) (0.90441,0.79441) (0.77384,0.75616) (0.90932,0.74419) (0.76607,0.74991) (0.84100,0.82669) (0.72223,0.69249) (0.85461,0.78063) (0.81206,0.65929) (0.72648,0.73876) (0.63959,0.61426) (0.80413,0.80050) (0.69950,0.74672) (0.78296,0.73935) (0.70276,0.70267) (0.87707,0.86354) (0.91622,0.79570) (0.96027,0.79787) (0.80585,0.78590) (0.77018,0.76059) (0.81372,0.79084) (0.88963,0.86683) (0.87950,0.83954) (0.87368,0.83389) (0.75885,0.75382) (0.65223,0.61726) (0.83946,0.77236) (0.65179,0.67576) (0.84010,0.82535) (0.78485,0.75757) (0.74455,0.71806) (0.74341,0.70407) (0.76307,0.71106) (0.80130,0.79016) (0.86440,0.84891) (0.67262,0.75955) (0.77538,0.77245) (0.82222,0.87481) (0.76653,0.73348) (0.72877,0.71012) (0.76004,0.74685) (0.75181,0.70336) (0.83289,0.77030) (0.78132,0.76907) (0.74148,0.74364) (0.64481,0.67523) (0.85471,0.85110) (0.84419,0.82932) (0.76748,0.75360) (0.66622,0.72145) (0.76186,0.73488) (0.75690,0.74757) (0.89072,0.88061) (0.82900,0.81400) (0.72541,0.67367) (0.83158,0.83090) (0.91507,0.86390) (0.84504,0.71075) (0.98888,0.78529) (0.81338,0.76173) (0.94582,0.71469) (0.92700,0.81194) (0.93426,0.74075) (0.95187,0.87273) (0.87999,0.73061) (0.91318,0.83304) (0.75517,0.72200) (0.80923,0.65597) (0.86547,0.72009) (0.91214,0.86378) (0.94262,0.87700) (0.87746,0.81565) (0.75234,0.70200) (0.89621,0.81584) (0.83017,0.76067) (0.84406,0.71097) (0.75516,0.68902) (0.94164,0.71440) (0.90257,0.77277) (0.93696,0.76132) (0.83397,0.71052) (0.93787,0.77398) (0.95530,0.78414) (0.94340,0.86661) (0.82934,0.79463) (0.84995,0.82641) (0.72547,0.72614) (0.82860,0.81765) (0.75124,0.72947) (0.80609,0.75079) (0.74374,0.69960) (0.97196,0.89320) (0.92317,0.73906) (0.95193,0.79084)};
\addplot+[black!65,densely dashed,very thick,no marks,domain=0.5:1.0] {x};
\end{groupplot}
\end{tikzpicture}
\caption{Human AUROC compared separately for each tissue--\plainHLA{} combination (task-wise AUROC) under \plainMatchedOOF{} versus \plainPeptideOOF{}. The gray dashed equal-performance line is the reference $y=x$, not an experimental series: points below it have lower AUROC after exact peptide separation across all tasks.}
\label{fig:human-standard-strict-scatter}
\end{figure}

Human AUROC decreases by 0.06275. The decrease is widespread rather than confined to a few unusual tissues: \plainPeptideOOF{} AUROC is lower for 145 of 157 human tissue--\plainHLA{} combinations (tasks). The mean change is -0.0628, with a 95\% interval obtained by resampling complete biological combinations (task-resampling interval) of -0.0728 to -0.0534. \plainPairAcc{} decreases by 0.0529. Full paired statistics are reported in Table~\ref{tab:paired-stats}. Exact peptide reuse therefore explains a meaningful part of the easier performance estimate, but it does not explain all signal because the harder-test AUROC remains well above random.

\subsubsection{Mouse benchmark}

The same protocol question is then evaluated in mouse to test whether the peptide-separation gap recurs in a separately constructed tissue--\plainHtwo{} benchmark; it is not a direct comparison of human and mouse model quality. Both mouse analyses evaluate the same 6,766 pairs. Across the 72 task-by-fold cells, 48 have exactly the same held-out pair count and the remaining cells differ by one pair, giving a mean absolute difference of 0.333 pair. In the ordinary folds, 71.43--72.12\% of unique held-out peptides reappear in fitting, compared with 0\% after global peptide separation.

\begin{table}[tbp]
\centering
\caption{Mouse AUROC under \plainMatchedOOF{} and \plainPeptideOOF{}. Difference is peptide-separated minus ordinary same-pool cross-validation; a negative value denotes lower performance after exact peptide separation.}
\label{tab:mouse-standard-strict-summary}
\small
\begin{tabularx}{\textwidth}{>{\raggedright\arraybackslash}X
                            >{\raggedright\arraybackslash}Xr}
\toprule
\plainMatchedOOF{} & \plainPeptideOOF{} & Difference \\ 
\midrule
0.8392 & 0.7529 & -0.0863 \\ 
\bottomrule
\end{tabularx}
\end{table}

The relationship across mouse tissue--\plainHtwo{} combinations is shown separately in Figure~\ref{fig:mouse-standard-strict-scatter}.
\begin{figure}[tbp]
\centering
\begin{tikzpicture}
\begin{groupplot}[
    group style={group size=1 by 1},
    width=0.72\linewidth,
    height=0.52\linewidth,
    xmin=0.5,xmax=1.0,
    ymin=0.5,ymax=1.0,
    xlabel={Same-pool ordinary cross-validation (matched standard OOF) AUROC},
    ylabel={Peptide-separated cross-validation (global peptide-disjoint OOF) AUROC},
    grid=major,
    grid style={gray!20}
]
\nextgroupplot[title={Mouse tissue--H2 combinations}]
\addplot+[only marks,mark=*,mark size=1.8pt,orange!85!black] coordinates {(0.97154,0.79471) (0.72771,0.61324) (0.69557,0.70072) (0.98073,0.83055) (0.82853,0.67053) (0.70739,0.66897) (0.82058,0.77812) (0.88098,0.82681) (0.96388,0.82430) (0.92562,0.82334) (0.66658,0.65301) (0.71439,0.58586) (0.88090,0.72499) (0.77186,0.76575) (0.86048,0.82001) (0.74906,0.69690) (0.95939,0.83240) (0.86177,0.75903) (0.75906,0.72217) (0.83286,0.78785) (0.85760,0.79344) (0.95193,0.80630) (0.80460,0.80518) (0.96826,0.78619)};
\addplot+[black!65,densely dashed,very thick,no marks,domain=0.5:1.0] {x};
\end{groupplot}
\end{tikzpicture}
\caption{Mouse AUROC compared separately for each tissue--\plainHtwo{} combination (task-wise AUROC) under \plainMatchedOOF{} versus \plainPeptideOOF{}. The gray dashed equal-performance line is the reference $y=x$, not an experimental series: points below it have lower AUROC after exact peptide separation across all tasks.}
\label{fig:mouse-standard-strict-scatter}
\end{figure}

Mouse AUROC decreases by 0.08629, and \plainPeptideOOF{} AUROC is lower for 22 of 24 tissue--\plainHtwo{} combinations (tasks). The mean change is -0.0863, with a task-resampling 95\% interval of -0.1106 to -0.0632. \plainPairAcc{} decreases by 0.0615. The numerically larger mouse decrease should not be interpreted as a species difference because the datasets, task inventories, component structures, and models were not otherwise matched.

For both species, we do not subtract the separate \plainFixedTest{} results from the \plainPeptideOOF{} results because those evaluations contain different examples.

\plainPeptideOOF{} removes reuse of exact peptide identities while retaining the benchmark's represented tasks and other shared context; the complete scope restrictions are stated in Table~\ref{tab:evaluation-protocols} and the Limitations.

\subsection{Research question 5 (RQ5): Replication in Tissue--Mouse-antigen-presenting-system Tasks (Tissue--H2 Tasks)}

The mouse benchmark asks whether the same biological question, controlling the presenting restriction and parent protein in each peptide pair (matched biological question), is learnable with the \plainHtwo{}. It uses a separate model suited to the smaller set of 24 tissue--\plainHtwo{} combinations; no human parameters are transferred. Table~\ref{tab:mouse-baselines} compares several mouse models before the \plainFixedTest{} is examined.

\begin{table}[tbp]
\centering
\caption{Mouse development trajectory under \plainPairOOF{} across 24 tissue--\plainHtwo{} combinations (tasks). The three middle rows are arithmetic means of three separately trained model summaries; the final row averages row-level predictions across five \plainSeeds{} (five-seed ensemble). This table is not a comparison in which all models use the same conditions (matched architecture comparison); the identical-fold \plainPeptideOOF{} comparison appears in Table~\ref{tab:mouse-strict-architectures}.}
\label{tab:mouse-baselines}
\footnotesize
\setlength{\tabcolsep}{3pt}
\begin{tabularx}{\textwidth}{>{\raggedright\arraybackslash}Xrrr
                            >{\raggedright\arraybackslash}p{1.8cm}}
\toprule
Model & AUROC & AUPRC & Worst-6 & Evidence \\ 
\midrule
\plainBLOSUM{} random forest & 0.7530 & 0.7292 & -- & reference \\ 
One shared peptide representation (shared encoder; mean across three \plainSeeds{}) & 0.8073 & 0.7929 & 0.6617 & member summary \\ 
\plainMMoE{} (mean across three \plainSeeds{}) & 0.8148 & 0.8043 & 0.6771 & member summary \\ 
Representation adjustment for the mouse antigen-presenting-system restriction H2-Kk (H2-Kk adapter; mean across three \plainSeeds{}) & 0.8180 & 0.8065 & 0.6816 & member summary \\ 
\plainMMoE{} (predictions averaged across five \plainSeeds{}; ensemble) & \textbf{0.8392} & \textbf{0.8316} & \textbf{0.7101} & row-level ensemble \\ 
\bottomrule
\end{tabularx}
\end{table}

The mouse model design, training settings, repeated runs, and averaging rule were \plainFrozen{} first. The selected model was then trained on all available mouse training pairs and evaluated once on the internal \plainFixedTest{}. Table~\ref{tab:mouse-protocol-summary} places this internal confirmation beside the harder result for peptides absent from all training tasks. The data protocol and number of averaged \plainSeeds{} are kept in separate columns.

\begin{table}[tbp]
\centering
\caption{Complete mouse evaluation summary. The \plainFixedTest{} is an internal confirmation in previously represented tasks with entity overlap across tasks; \plainPeptideOOF{} evaluates exact peptide identities absent from training in those same represented tasks.}
\label{tab:mouse-protocol-summary}
\footnotesize
\setlength{\tabcolsep}{2pt}
\begin{tabularx}{\textwidth}{>{\raggedright\arraybackslash}X
                            >{\raggedright\arraybackslash}Xrrrr}
\toprule
Evaluation protocol & Averaged models (ensemble) & AUROC & AUPRC & Worst-6 & \shortstack{Within-pair ordering\\accuracy (PairAcc)} \\ 
\midrule
\plainMatchedOOF{} & five \plainSeeds{} & 0.8392 & 0.8316 & 0.7101 & 0.8211 \\ 
\plainFixedTest{} & five \plainSeeds{}; \plainFrozen{} & \textbf{0.8562} & \textbf{0.8506} & \textbf{0.7245} & \textbf{0.8425} \\ 
\plainPeptideOOF{} & five \plainSeeds{} & 0.7529 & 0.7322 & 0.6481 & 0.7596 \\ 
\bottomrule
\end{tabularx}
\end{table}

Relative to \plainMatchedOOF{}, the \plainFixedTest{} values are descriptively higher by 0.0170 AUROC, 0.0190 AUPRC, and 0.0144 worst-6 AUROC. Eighteen of 24 tasks improve in AUROC and six decline. This is not an estimate of the effect of peptide separation because the evaluated examples differ. It is an internal confirmation rather than independent validation because the \plainFixedTest{} contains substantial entity overlap between fitting and test data.

\begin{table}[tbp]
\centering
\caption{Mouse \plainFixedTest{} AUROC grouped by \plainHtwo{} restriction. These are descriptive subgroup means, not controlled effects of the \plainHtwo{} restriction.}
\label{tab:mouse-h2-groups}
\small
\begin{tabularx}{\textwidth}{>{\raggedright\arraybackslash}X
                            >{\centering\arraybackslash}X
                            >{\centering\arraybackslash}X
                            >{\centering\arraybackslash}X
                            >{\centering\arraybackslash}X}
\toprule
\plainHtwo{} restriction & Mouse Db restriction (H2-Db) & Mouse Kb restriction (H2-Kb) & Mouse Kd restriction (H2-Kd) & Mouse Kk restriction (H2-Kk) \\ 
\midrule
AUROC averaged across tissue combinations (mean task AUROC) & 0.9152 & 0.7404 & 0.8397 & 0.8024 \\ 
\bottomrule
\end{tabularx}
\end{table}

Equal-weight summaries within each tissue also reveal marked heterogeneity. Table~\ref{tab:mouse-tissue-extremes} lists the four lowest and four highest tissue-level means on the \plainFixedTest{} (tissue-macro AUROCs); tissues represented by one task should be interpreted especially cautiously.

\begin{table}[tbp]
\centering
\caption{Descriptive mouse tissue-level extremes on the \plainFixedTest{}. Means are taken across the available \plainHtwo{} tasks within each tissue.}
\label{tab:mouse-tissue-extremes}
\footnotesize
\setlength{\tabcolsep}{2pt}
\begin{tabular}{lrrrlrrr}
\toprule
Lower tissue & Tasks & AUROC & AUPRC & Higher tissue & Tasks & AUROC & AUPRC \\ 
\midrule
Colon & 2 & 0.7059 & 0.7186 & Lung & 1 & 0.9641 & 0.9579 \\ 
Pancreas & 1 & 0.7091 & 0.6737 & Bone marrow & 1 & 0.9734 & 0.9558 \\ 
Prostate & 1 & 0.7974 & 0.7719 & Small intestine & 1 & 0.9793 & 0.9748 \\ 
Skin & 4 & 0.8180 & 0.8075 & Kidney & 1 & 0.9837 & 0.9844 \\ 
\bottomrule
\end{tabular}
\end{table}

When every test peptide is absent from mouse training, AUROC is 0.7529 (95\% interval obtained by resampling complete biological combinations [task-resampling interval], 0.7230--0.7799), AUPRC is 0.7322, and 75.96\% of \plainMatchedPairs{} are ordered correctly. Together with the human result, this shows that the biological comparison remains learnable in both species but is consistently harder for new peptide identities. It does not show that the selected mouse architecture is uniquely responsible for the result.

\FloatBarrier
\subsection{Research question 6 (RQ6): How Much Signal Remains without Tissue Identity?}

Removing tissue identity while retaining peptide and \plainMHC{} information (\plainMHCOnly{}) lowers performance in every evaluation. Under \plainPeptideOOF{}, human AUROC falls from 0.7652 for the full tissue-conditioned model to 0.6630 for the \plainMHCOnly{}, and mouse AUROC falls from 0.7529 to 0.6412. The decrease occurs in 156 of 157 human tasks and all 24 mouse tasks. Thus, general peptide--antigen-presenting-molecule compatibility (peptide--MHC compatibility) alone does not reproduce the full ranking signal. This is evidence that tissue context matters to the benchmark, but it is not proof that the model has learned a causal tissue-processing mechanism.

\begin{table}[tbp]
\centering
\caption{Comparison of three methods that do not use tissue identity. The \plainMHCOnly{} has approximately the same number of model parameters as the corresponding full model (capacity-matched control). MHCflurry uses its predicted-presentation score, and NetMHCpan uses its \plainEL{}. External scores are \plainFrozen{}; therefore their values are identical under \plainMatchedOOF{} and \plainPeptideOOF{}, which use the same complete training pool. The best value among the three controls in each row is bold.}
\label{tab:rq6-tissue-blind-controls}
\scriptsize
\setlength{\tabcolsep}{4pt}
\begin{tabularx}{\textwidth}{@{}>{\raggedright\arraybackslash}Xrrr@{}}
\toprule
Evaluation metric & \plainMHCOnly{} & \shortstack{MHCflurry predicted\\presentation} & \shortstack{NetMHCpan\\\plainEL{}} \\
\midrule
\multicolumn{4}{@{}l}{\textit{Human}} \\
\plainFixedTest{}: AUROC & \textbf{0.7449} & 0.6851 & 0.6833 \\
\plainFixedTest{}: AUPRC & \textbf{0.7563} & 0.6783 & 0.6686 \\
\plainFixedTest{}: \plainPairAcc{} & \textbf{0.7243} & 0.6886 & 0.6897 \\
\addlinespace
\plainMatchedOOF{}: AUROC & \textbf{0.7353} & 0.6871 & 0.6823 \\
\plainMatchedOOF{}: AUPRC & \textbf{0.7379} & 0.6751 & 0.6632 \\
\plainMatchedOOF{}: \plainPairAcc{} & \textbf{0.7144} & 0.6898 & 0.6875 \\
\addlinespace
\plainPeptideOOF{}: AUROC & 0.6630 & \textbf{0.6871} & 0.6823 \\
\plainPeptideOOF{}: AUPRC & 0.6502 & \textbf{0.6751} & 0.6632 \\
\plainPeptideOOF{}: \plainPairAcc{} & 0.6691 & \textbf{0.6898} & 0.6875 \\
\midrule
\multicolumn{4}{@{}l}{\textit{Mouse}} \\
\plainFixedTest{}: AUROC & 0.6672 & 0.6564 & \textbf{0.6893} \\
\plainFixedTest{}: AUPRC & \textbf{0.7087} & 0.6436 & 0.6926 \\
\plainFixedTest{}: \plainPairAcc{} & 0.6354 & 0.6733 & \textbf{0.6917} \\
\addlinespace
\plainMatchedOOF{}: AUROC & \textbf{0.6819} & 0.6474 & 0.6802 \\
\plainMatchedOOF{}: AUPRC & \textbf{0.7008} & 0.6263 & 0.6785 \\
\plainMatchedOOF{}: \plainPairAcc{} & 0.6540 & 0.6565 & \textbf{0.6888} \\
\addlinespace
\plainPeptideOOF{}: AUROC & 0.6412 & 0.6474 & \textbf{0.6802} \\
\plainPeptideOOF{}: AUPRC & 0.6378 & 0.6263 & \textbf{0.6785} \\
\plainPeptideOOF{}: \plainPairAcc{} & 0.6418 & 0.6565 & \textbf{0.6888} \\
\bottomrule
\end{tabularx}
\end{table}

MHCflurry and NetMHCpan perform better than random on the \plainMatchedPairs{} (Table~\ref{tab:rq6-tissue-blind-controls}), with AUROCs of 0.6564--0.6893 on the \plainFixedTest{}. General binding and presentation propensity therefore contributes useful information. However, these scores remain well below the tissue-conditioned models, leaving a substantial part of the ordering associated with tissue context or with other structure captured by the benchmark.

The NetMHCpan \plainEL{} is consistently at least as strong as its \plainBA{}. The AUROC difference between \plainEL{} and \plainBA{} is 0.0256 on the human \plainFixedTest{} and 0.0264 on the human complete training pool, compared with 0.0046 and 0.0024 in mouse. We therefore use the \plainEL{} as the primary NetMHCpan control and report the complete \plainBA{} sensitivity analysis in Table~\ref{tab:appendix-netmhcpan-el-ba}.

The tissue-conditioned models remain higher than the strongest \plainGeneralPMHC{} under every protocol. Relative to the three tissue-blind controls in Table~\ref{tab:rq6-tissue-blind-controls}, the relevant strongest external control is MHCflurry predicted presentation in human and NetMHCpan \plainEL{} in mouse. The tissue-conditioned-versus-external AUROC values are 0.8448 versus 0.6851 on the human \plainFixedTest{}, 0.8279 versus 0.6871 under human \plainMatchedOOF{}, and 0.7652 versus 0.6871 under human \plainPeptideOOF{}; the corresponding mouse values are 0.8562 versus 0.6893, 0.8392 versus 0.6802, and 0.7529 versus 0.6802. The resulting advantages on the \plainFixedTest{} and under \plainMatchedOOF{} are 0.1409--0.1669 AUROC. Under \plainPeptideOOF{}, the advantages narrow to 0.0781 in human and 0.0727 in mouse, showing that peptide separation weakens but does not erase the equal-weight average across tissue--antigen-presenting-molecule combinations (macro-average) gap. In the analysis comparing the same biological combinations under both models (paired task analysis), the human advantages under all three protocols and the mouse advantages on the \plainFixedTest{} and under \plainMatchedOOF{} pass correction for testing several hypotheses (BH-FDR correction). The mouse \plainPeptideOOF{} comparison with the NetMHCpan \plainEL{} is heterogeneous: its mean difference is 0.0727 with a 95\% interval obtained by resampling complete biological combinations (task-resampling interval) of 0.0028--0.1579, but the multiple-testing-adjusted Wilcoxon value is \(q=0.6840\), with the tissue-conditioned model higher in 11 combinations (tasks) and lower in 13. Thus, this mouse result supports a positive equal-weight average across tissue--antigen-presenting-molecule combinations (macro-average), not a claim of consistent superiority across most combinations (tasks).

Learning a score-combination rule on other folds before applying it to each held-out fold (cross-fitted stacking) yields only marginal improvement over the tissue-conditioned score alone. The largest change is obtained for human \plainPeptideOOF{} with MHCflurry, from 0.7652 to 0.7693 AUROC (+0.0041). For the strongest species-specific external controls, the \plainFixedTest{} changes are 0.8448 to 0.8453 in human and 0.8562 to 0.8564 in mouse, while the mouse \plainPeptideOOF{} change is 0.7529 to 0.7534. The complete score-combination results (stacker results) appear in Table~\ref{tab:appendix-general-stacking}. General binding or presentation propensity therefore explains part of the benchmark signal, but neither the standalone controls nor their residual incremental contribution accounts for the full performance of the tissue-conditioned models.

\subsection{Research question 7 (RQ7): Effects of Model Structure under Peptide Separation}

To compare model structures (architectures) fairly, every model is retrained on exactly the same folds from \plainPeptideOOF{}. All human models retain useful signal (Table~\ref{tab:human-strict-architectures}). TissuePMHC is clearly better than simple sequence baselines and is modestly better than models that either use \plainSharedHeads{} or use the \plainPlainDual{} without extra tissue and \plainHLA{} training guidance.

\begin{table}[tbp]
\centering
\caption{Human comparison of model structures under \plainPeptideOOF{} across 157 tasks on identical folds built from transitively linked peptide pairs (connected-component folds). Methods with several \plainSeeds{} average row-level predictions; logistic regression has no random training variation and therefore gives identical predictions for the recorded \plainSeeds{}.}
\label{tab:human-strict-architectures}
\scriptsize
\setlength{\tabcolsep}{2pt}
\begin{tabularx}{\textwidth}{>{\raggedright\arraybackslash}Xrrrrrr}
\toprule
Model & Mean AUROC & Median & AUPRC & \shortstack{Within-pair ordering\\accuracy (PairAcc)} & Worst-10 & Worst task \\
\midrule
One-hot logistic regression & 0.7127 & 0.7069 & 0.6906 & 0.7157 & 0.5873 & 0.5533 \\
\plainBLOSUM{} random forest & 0.7212 & 0.7208 & 0.7007 & 0.7263 & 0.5895 & 0.5634 \\
\plainSharedHeads{} & 0.7561 & 0.7565 & 0.7340 & 0.7631 & 0.6248 & 0.6003 \\
\plainPlainDual{} using an \plainMLP{} & 0.7555 & 0.7577 & 0.7339 & 0.7575 & 0.6163 & 0.5905 \\
\plainAuxDual{} & 0.7642 & \textbf{0.7660} & 0.7433 & 0.7686 & 0.6283 & 0.6034 \\
TissuePMHC & \textbf{0.7652} & 0.7606 & \textbf{0.7452} & \textbf{0.7700} & \textbf{0.6382} & \textbf{0.6142} \\
\bottomrule
\end{tabularx}
\end{table}

Comparison with the strongest simpler human control yields a more cautious interpretation. This \plainAuxDual{} uses the same two-pathway structure and additional tissue and \plainHLA{} guidance during training (auxiliary supervision), but a simpler peptide representation. TissuePMHC is only 0.0010 AUROC higher and has 76 positive, one tied, and 80 negative task differences (win/tie/loss count); the uncertainty interval includes zero. In contrast, adding tissue and \plainHLA{} training guidance to the \plainPlainDual{} improves AUROC by 0.0087 and benefits 144 of 157 tasks. The clearest retained effect is therefore additional training guidance based on biological groups (auxiliary biological-group supervision), rather than the more elaborate peptide encoder alone.

The three mouse models are also very close under \plainPeptideOOF{} (Table~\ref{tab:mouse-strict-architectures}). The selected \plainMMoE{} differs from \plainSharedHeads{} by -0.0008 AUROC and from the \plainBLOSUM{} random forest by +0.0023; neither difference is consistent across tasks. The mouse experiment therefore supports learnability of the tissue--\plainHtwo{} comparison, not superiority of one complex model structure (architecture).

\begin{table}[tbp]
\centering
\caption{Mouse comparison of model structures under \plainPeptideOOF{} across 24 tasks. All entries average predictions across five \plainSeeds{} (five-seed ensembles) on identical folds built from transitively linked peptide pairs (connected-component folds).}
\label{tab:mouse-strict-architectures}
\scriptsize
\setlength{\tabcolsep}{1pt}
\begin{tabularx}{\textwidth}{>{\raggedright\arraybackslash}Xrrrrrr}
\toprule
Model & Mean AUROC & Median & AUPRC & \shortstack{Within-pair ordering\\accuracy (PairAcc)} & Worst-6 & Worst task \\
\midrule
\plainBLOSUM{} random forest & 0.7507 & 0.7768 & 0.7295 & 0.7573 & 0.6418 & 0.5819 \\
\plainSharedHeads{} & \textbf{0.7538} & 0.7786 & 0.7321 & 0.7567 & 0.6430 & 0.5477 \\
\plainMMoE{} & 0.7529 & \textbf{0.7822} & \textbf{0.7322} & \textbf{0.7596} & \textbf{0.6481} & \textbf{0.5859} \\
\bottomrule
\end{tabularx}
\end{table}

Averaging predictions across \plainSeeds{} contributes materially to the final neural-model averages (ensembles). In human, the AUROC gains from the averaged prediction relative to the mean of the separate runs are +0.0248 for \plainSharedHeads{}, +0.0161 for the \plainPlainDual{}, +0.0139 for the \plainAuxDual{}, and +0.0147 for TissuePMHC. In mouse, the corresponding gains are +0.0211 for \plainSharedHeads{} and +0.0318 for \plainMMoE{}, compared with +0.0016 for the random forest. The mean of the separate \plainMMoE{} runs is only 0.7211, below \plainSharedHeads{} (0.7327) and the random forest (0.7491). Its competitive final score under \plainPeptideOOF{} is therefore driven primarily by averaging five \plainSeeds{}, rather than by a consistent advantage of the individual model structure. Full paired summaries and statistics across \plainSeeds{} are reported in Tables~\ref{tab:appendix-strict-paired} and~\ref{tab:appendix-strict-seeds}.
